\documentclass[10pt,a4paper]{article}

\usepackage[utf8]{inputenc}
\usepackage[T1]{fontenc}
\usepackage[italian]{babel}
\usepackage{amsmath,amssymb}
\usepackage{geometry}
\usepackage{tabularx}
\usepackage{enumitem}

\geometry{
	left=1.4cm,
	right=1.4cm,
	top=1.1cm,
	bottom=1.1cm
}

\setlength{\parindent}{0pt}
\setlength{\parskip}{3pt}
\renewcommand{\arraystretch}{1.2}

\setlist[enumerate]{
	leftmargin=6mm,
	itemsep=1pt,
	topsep=2pt
}

% ============================================================
% COMANDI COMUNI
% ============================================================

\newcommand{\ruoli}{
\begin{center}
\begin{tabularx}{\textwidth}{|X|X|}
\hline
\textbf{Nome} & \textbf{Ruolo} \\
\hline
& Coordinatore/coordinatrice \\
\hline
& Controllore/controllatrice dei dati \\
\hline
& Segretario/a \\
\hline
& Portavoce \\
\hline
\end{tabularx}
\end{center}

\textit{I ruoli servono per organizzare il confronto di gruppo.
Le costruzioni devono invece essere svolte individualmente da tutti.}
}

\newcommand{\regole}{
\begin{center}
\fbox{
\begin{minipage}{0.94\textwidth}
\textbf{Regole dell'esperimento}

\begin{enumerate}
	\item Ognuno riceve \textbf{un foglio bianco} e lo taglia in due metà:
	una per il \textbf{TENTATIVO 1} e una per il \textbf{TENTATIVO 2}.

	\item Ognuno deve eseguire la costruzione \textbf{da solo}.

	\item Esegui la stessa costruzione due volte, una su ciascuna metà.

	\item Il secondo tentativo deve essere una nuova costruzione:
	\textbf{non ricalcare e non copiare il primo}.

	\item Durante le costruzioni \textbf{non guardare i disegni degli altri}
	e non suggerire loro come procedere.

	\item Non confrontate i risultati finché tutti non hanno terminato
	entrambi i tentativi.

	\item Usa un tratto abbastanza visibile da poter osservare
	il disegno anche attraverso il foglio.

	\item Per confrontare due triangoli, sovrapponete le metà dei fogli
	e guardatele \textbf{controluce}.

	\item Potete ruotare o ribaltare una delle due metà:
	l'orientamento del triangolo non conta.
\end{enumerate}
\end{minipage}
}
\end{center}
}

\newcommand{\autocontrollo}{
\subsection*{2. Controlla prima il tuo lavoro}

Prima di guardare i disegni degli altri, sovrapponi le due metà
del tuo foglio e osserva i tuoi due tentativi controluce.

Puoi ruotare o ribaltare una delle due metà.

I due triangoli riescono a coincidere?

\[
\square\ \text{Sì}
\qquad\qquad
\square\ \text{No}
\]

Se non coincidono perfettamente, prova a capire quale dato
potresti non aver riprodotto con sufficiente precisione.

\textbf{Non cancellare né rifare i tentativi: anche gli errori
fanno parte dell'esperimento.}
}

\newcommand{\righe}{
\vspace{1mm}
\hrulefill

\vspace{3mm}
\hrulefill

\vspace{3mm}
\hrulefill
}

\newcommand{\classificazione}{
\subsection*{5. Osservate e classificate}

Rispetto ai \textbf{lati}, il triangolo è:

\[
\boxed{\hspace{5.5cm}}
\]

Rispetto agli \textbf{angoli}, il triangolo è:

\[
\boxed{\hspace{5.5cm}}
\]

Quali caratteristiche avete osservato per deciderlo?

\righe
}

% ============================================================
% DOCUMENTO
% ============================================================

\begin{document}

% ============================================================
% GRUPPO 1
% ============================================================

\begin{center}
{\Large\bfseries GRUPPO 1}\\[1mm]
{\large Un triangolo è davvero ``bloccato''?}
\end{center}

\ruoli
\regole

\subsection*{I vostri dati}

Costruite un triangolo \(ABC\) sapendo che

\[
AB=5\text{ cm},
\qquad
AC=5\text{ cm},
\qquad
\widehat{BAC}=60^\circ.
\]

\subsection*{1. Costruzione individuale}

Ciascuno costruisca \textbf{due volte} un triangolo che rispetti
esattamente i dati assegnati, una volta su ciascuna metà del proprio foglio.

\autocontrollo

\subsection*{3. Ora confrontatevi}

Solo quando tutti hanno terminato, confrontate i triangoli costruiti
dai diversi componenti del gruppo.

Sovrapponete a coppie le metà dei fogli e osservatele controluce.
Ruotatele o ribaltatele se necessario.

I triangoli costruiti dai diversi componenti riescono a coincidere?

\[
\square\ \text{Sì}
\qquad\qquad
\square\ \text{No}
\]

Secondo voi potrebbe esistere un triangolo \textbf{non congruente}
che rispetti comunque esattamente gli stessi dati?

\[
\square\ \text{Sì}
\qquad\qquad
\square\ \text{No}
\]

\subsection*{4. La conclusione del gruppo}

Questi dati sono sufficienti per ``bloccare'' il triangolo?
Spiegate perché.

\righe

\classificazione

\newpage

% ============================================================
% GRUPPO 2
% ============================================================

\begin{center}
{\Large\bfseries GRUPPO 2}\\[1mm]
{\large Un triangolo è davvero ``bloccato''?}
\end{center}

\ruoli
\regole

\subsection*{I vostri dati}

Costruite un triangolo \(ABC\) sapendo che

\[
AB=6\text{ cm},
\qquad
\widehat{BAC}=45^\circ,
\qquad
\widehat{ABC}=45^\circ.
\]

\subsection*{1. Costruzione individuale}

Ciascuno costruisca \textbf{due volte} un triangolo che rispetti
esattamente i dati assegnati, una volta su ciascuna metà del proprio foglio.

\autocontrollo

\subsection*{3. Ora confrontatevi}

Solo quando tutti hanno terminato, confrontate i triangoli costruiti
dai diversi componenti.

Sovrapponete a coppie le metà dei fogli e osservatele controluce.

I triangoli riescono a coincidere?

\[
\square\ \text{Sì}
\qquad\qquad
\square\ \text{No}
\]

Potrebbe esistere un triangolo diverso per forma o dimensioni
che rispetti comunque tutti i dati assegnati?

\[
\square\ \text{Sì}
\qquad\qquad
\square\ \text{No}
\]

\subsection*{4. La conclusione del gruppo}

Questi dati sono sufficienti per ``bloccare'' il triangolo?
Spiegate perché.

\righe

\classificazione

\newpage

% ============================================================
% GRUPPO 3
% ============================================================

\begin{center}
{\Large\bfseries GRUPPO 3}\\[1mm]
{\large Un triangolo è davvero ``bloccato''?}
\end{center}

\ruoli
\regole

\subsection*{I vostri dati}

Costruite un triangolo \(DEF\) sapendo che

\[
DE=5\text{ cm},
\qquad
\widehat{EDF}=30^\circ,
\qquad
\widehat{DEF}=60^\circ.
\]

\subsection*{1. Costruzione individuale}

Ciascuno costruisca \textbf{due volte} un triangolo che rispetti
esattamente i dati assegnati, una volta su ciascuna metà del proprio foglio.

\autocontrollo

\subsection*{3. Ora confrontatevi}

Solo quando tutti hanno terminato, confrontate i triangoli costruiti
dai diversi componenti.

Sovrapponete a coppie le metà dei fogli e osservatele controluce.

I triangoli riescono a coincidere?

\[
\square\ \text{Sì}
\qquad\qquad
\square\ \text{No}
\]

Potrebbe esistere un triangolo diverso per forma o dimensioni
che rispetti comunque tutti i dati assegnati?

\[
\square\ \text{Sì}
\qquad\qquad
\square\ \text{No}
\]

\subsection*{4. La conclusione del gruppo}

Questi dati sono sufficienti per ``bloccare'' il triangolo?
Spiegate perché.

\righe

\classificazione

\newpage

% ============================================================
% GRUPPO 4
% ============================================================

\begin{center}
{\Large\bfseries GRUPPO 4}\\[1mm]
{\large Un triangolo è davvero ``bloccato''?}
\end{center}

\ruoli
\regole

\subsection*{I vostri dati}

Costruite un triangolo \(ABC\) i cui lati misurino

\[
AB=3\text{ cm},
\qquad
AC=4\text{ cm},
\qquad
BC=5\text{ cm}.
\]

\subsection*{1. Costruzione individuale}

Ciascuno costruisca \textbf{due volte} un triangolo che rispetti
esattamente le tre lunghezze assegnate, una volta su ciascuna metà
del proprio foglio.

\autocontrollo

\subsection*{3. Ora confrontatevi}

Solo quando tutti hanno terminato, confrontate i triangoli costruiti
dai diversi componenti.

Sovrapponete a coppie le metà dei fogli e osservatele controluce.

I triangoli riescono a coincidere?

\[
\square\ \text{Sì}
\qquad\qquad
\square\ \text{No}
\]

Secondo voi è possibile spostare uno dei vertici, mantenendo
\textbf{esattamente} le tre lunghezze assegnate,
e ottenere un triangolo di forma diversa?

\[
\square\ \text{Sì}
\qquad\qquad
\square\ \text{No}
\]

\subsection*{4. La conclusione del gruppo}

Conoscere esattamente le lunghezze dei tre lati è sufficiente
per ``bloccare'' il triangolo?

\righe

\classificazione

\newpage

% ============================================================
% GRUPPO 5
% ============================================================

\begin{center}
{\Large\bfseries GRUPPO 5}\\[1mm]
{\large Un triangolo è davvero ``bloccato''?}
\end{center}

\ruoli
\regole

\subsection*{I vostri dati}

Costruite un triangolo \(ABC\) sapendo che

\[
\widehat A=30^\circ,
\qquad
\widehat B=60^\circ,
\qquad
\widehat C=90^\circ.
\]

\textbf{Non è assegnata la lunghezza di nessun lato.}

\subsection*{1. Costruzione individuale}

Prima di iniziare, ciascuno scelga autonomamente una lunghezza
per il lato \(AB\).

\textbf{Non dite agli altri quale lunghezza avete scelto.}

Usate la stessa lunghezza di \(AB\) nei vostri due tentativi,
così potrete verificare se siete riusciti a ripetere la stessa costruzione.

Ciascuno costruisca quindi il triangolo \textbf{due volte},
una volta su ciascuna metà del proprio foglio.

\autocontrollo

\subsection*{3. Ora confrontatevi}

Solo adesso confrontate le costruzioni dei diversi componenti
e rivelate quale lunghezza di \(AB\) avete scelto.

I triangoli hanno tutti la stessa forma?

\[
\square\ \text{Sì}
\qquad\qquad
\square\ \text{No}
\]

Hanno necessariamente la stessa dimensione?

\[
\square\ \text{Sì}
\qquad\qquad
\square\ \text{No}
\]

Sono necessariamente congruenti?

\[
\square\ \text{Sì}
\qquad\qquad
\square\ \text{No}
\]

\subsection*{4. La conclusione del gruppo}

Conoscere i tre angoli è sufficiente per ``bloccare'' un triangolo?

Che cosa rimane uguale e che cosa può invece cambiare?

\righe

\subsection*{5. Osservate e classificate}

I vostri triangoli possono avere dimensioni diverse.
Appartengono comunque alle stesse categorie?

Rispetto ai \textbf{lati} sono:

\[
\boxed{\hspace{5.5cm}}
\]

Rispetto agli \textbf{angoli} sono:

\[
\boxed{\hspace{5.5cm}}
\]

Che cosa hanno tutti in comune?

\righe

\newpage

% ============================================================
% GRUPPO 6
% ============================================================

\begin{center}
{\Large\bfseries GRUPPO 6}\\[1mm]
{\large Un triangolo è davvero ``bloccato''?}
\end{center}

\ruoli
\regole

\subsection*{I vostri dati}

Costruite un triangolo \(ABC\) sapendo che

\[
AB=6\text{ cm},
\qquad
AC=5\text{ cm},
\qquad
\widehat{ABC}=45^\circ.
\]

\textbf{Attenzione:} l'angolo assegnato non è compreso
tra i due lati di cui conoscete la lunghezza.

\subsection*{1. Costruzione individuale}

In ciascuno dei due tentativi:

\begin{enumerate}
	\item costruisci \(AB=6\) cm;
	\item da \(B\) costruisci una semiretta che formi un angolo
	di \(45^\circ\) con \(BA\);
	\item il punto \(C\) deve trovarsi su questa semiretta;
	\item deve inoltre essere
	\[
	AC=5\text{ cm};
	\]
	\item cerca con attenzione \textbf{tutti i possibili punti \(C\)}
	che rispettano contemporaneamente i dati.
\end{enumerate}

Ripeti poi l'intera costruzione sulla seconda metà del foglio,
senza copiare la prima.

\subsection*{2. Controlla prima il tuo lavoro}

Nei due tentativi hai trovato lo stesso numero di possibili triangoli?

\[
\square\ \text{Sì}
\qquad\qquad
\square\ \text{No}
\]

Quanti triangoli diversi sembrano rispettare esattamente i dati?

\[
\boxed{\hspace{2cm}}
\]

Sovrapponi anche i risultati corrispondenti dei tuoi due tentativi
per controllare se coincidono.

\subsection*{3. Ora confrontatevi}

Solo adesso osservate le costruzioni degli altri componenti del gruppo.

Avete trovato tutti lo stesso numero di possibili triangoli?

\[
\square\ \text{Sì}
\qquad\qquad
\square\ \text{No}
\]

I triangoli possibili sono tutti congruenti tra loro?

\[
\square\ \text{Sì}
\qquad\qquad
\square\ \text{No}
\]

\subsection*{4. La conclusione del gruppo}

Due lati e un angolo non compreso sono sempre sufficienti
per ``bloccare'' un triangolo?

\righe

\subsection*{5. Osservate e classificate}

Se avete trovato due triangoli diversi, classificateli separatamente.

\textbf{Primo triangolo}

Rispetto ai lati:

\[
\boxed{\hspace{5cm}}
\]

Rispetto agli angoli:

\[
\boxed{\hspace{5cm}}
\]

\textbf{Secondo triangolo}

Rispetto ai lati:

\[
\boxed{\hspace{5cm}}
\]

Rispetto agli angoli:

\[
\boxed{\hspace{5cm}}
\]

I due triangoli appartengono alla stessa categoria rispetto agli angoli?

\[
\square\ \text{Sì}
\qquad\qquad
\square\ \text{No}
\]

Che cosa avete osservato?

\righe

\end{document}
